\section{Results}
\label{sec:results}

\subsection{Repeats in agentic sessions and the n-gram cache}
\label{sec:repeats}

\begin{figure}[tb]
\centering
\begin{tikzpicture}
\begin{groupplot}[
  group style={group size=2 by 1, horizontal sep=2.9cm},
  height=5.4cm,
  axis lines=left, axis line style={mute}, tick style={mute},
  ymajorgrids, grid style={black!8},
  tick label style={font=\footnotesize}, label style={font=\footnotesize},
  title style={font=\footnotesize},
  xmode=log, xmin=0.8, xmax=330,
  xtick={1.024,4.096,16.384,65.536,262.144}, xticklabels={1k,4k,16k,64k,256k},
  xlabel={Position in the session (tokens)},
  ymin=0, ymax=100, ytick={0,25,50,75,100}, clip=false,
]
\nextgroupplot[title={(a) Preceding $N$-gram occurred earlier}, width=0.385\textwidth,
  ylabel={Share of tokens (\%)}]
\draw[black!25, densely dashed] (axis cs:16.384,0) -- (axis cs:16.384,100);
\addplot[ramp1!55, thick, mark=*, mark size=1.3pt] table[x=ktok, y=s1]{generated/gram_ctx.dat};
\addplot[ramp1, thick, mark=*, mark size=1.3pt] table[x=ktok, y=s2]{generated/gram_ctx.dat};
\addplot[ramp1!50!ramp2, thick, mark=*, mark size=1.3pt] table[x=ktok, y=s4]{generated/gram_ctx.dat};
\addplot[ramp2, thick, mark=*, mark size=1.3pt] table[x=ktok, y=s8]{generated/gram_ctx.dat};
\addplot[ramp2!50!ramp3, thick, mark=*, mark size=1.3pt] table[x=ktok, y=s16]{generated/gram_ctx.dat};
\addplot[ramp3, thick, mark=*, mark size=1.3pt] table[x=ktok, y=s32]{generated/gram_ctx.dat};
\node[font=\scriptsize, anchor=west, text=ramp1!80!ramp2] at (axis cs:280,102) {$N=1$};
\node[font=\scriptsize, anchor=west, text=ramp2] at (axis cs:280,90.5) {$N=8$};
\node[font=\scriptsize, anchor=west, text=ramp3] at (axis cs:280,77.5) {$N=32$};

\nextgroupplot[title={(b) Lookup accuracy and token kinds}, width=0.385\textwidth,
  ylabel={Share of tokens (\%)}]
\draw[black!25, densely dashed] (axis cs:16.384,0) -- (axis cs:16.384,100);
\addplot[cat1, thick, mark=*, mark size=1.6pt] table[x=ktok, y=lookup]{generated/gram_ctx.dat};
\addplot[cat3!85!black, thick, mark=triangle*, mark size=2pt] table[x=ktok, y=far]{generated/gram_ctx.dat};
\addplot[cat2, thick, mark=square*, mark size=1.5pt] table[x=ktok, y=older]{generated/gram_ctx.dat};
\node[font=\scriptsize, anchor=west, text=cat1!85!black] at (axis cs:280,91.6) {lookup right};
\node[font=\scriptsize, anchor=west, text=cat3!65!black] at (axis cs:280,22) {far 8-gram match};
\node[font=\scriptsize, anchor=west, text=cat2!85!black] at (axis cs:280,5.6) {older value};
\end{groupplot}
\end{tikzpicture}

\caption{Repeat statistics of the 704 held-out trajectories, read from their first token, by
position in the session. (a) Share of tokens whose preceding $N$-gram occurred earlier in the
session, for the six cue lengths of \ingram{}. (b) Share of tokens that the longest-match lookup
of Section~\ref{sec:repeats} predicts, share whose preceding 8-gram occurred earlier with its most
recent occurrence more than 512 tokens back, and share for which an older continuation of a 4-token
cue is right and the most recent one is wrong. Each point pools one range of positions (up to 1k,
1k to 2k, 2k to 4k, and so on); 566 trajectories reach 64k tokens and 134 reach 128k. The dashed
line marks 16,384 tokens, the context of our models. These are counts on token ids and involve no
model.}
\label{fig:repeats}
\end{figure}

In the 32 evaluation windows, the $N$-gram that precedes a token has already occurred in the same
window for 94.6\% of tokens at $N = 1$, 77.0\% at $N = 4$ and 65.5\% at $N = 8$. A lookup with no
parameters, which finds the longest earlier match ($N = 8$ down to $1$) and copies the token that
followed its most recent occurrence, predicts 78.8\% of next tokens. These windows are cut from
longer sessions: the held-out trajectories have a median of 96k tokens, and repeats become more
frequent and longer with depth (Figure~\ref{fig:repeats}). The preceding 8-gram has occurred
earlier for 53.2\% of tokens at positions 8k to 16k and for 90.2\% beyond 128k, and the lookup
predicts 73.5\% and 91.6\% of tokens at these depths. The contexts of 16,384 tokens that we train
and evaluate on therefore contain fewer of the repeats that \ingram{} reads than a full session
does.

Because a lookup predicts this share of tokens, we include a baseline that operates at inference
time, the n-gram cache. At each position it finds the longest $N \le 8$ whose $N$-gram occurred
earlier and predicts from the counts of its earlier continuations. We mix it with a model using
one weight per matched $N$, fit the weights on 16 windows and score the other 16, in both
directions. On seeds 0 and 1 the cache lowers the loss of GDN from 1.002 to 0.683 and the loss of
the 3:1 hybrid from 0.722 to 0.621. We therefore report each $\Delta$ twice: between the two
models, and between the two models with the cache mixed into both.

\subsection{Replacing the attention layers of a 3:1 hybrid}
\label{sec:replace}

\begin{figure}[tb]
\centering
\resizebox{\textwidth}{!}{\begin{tikzpicture}
\def\lx{-8.5cm}
\def\ours#1#2#3{\node[font=\footnotesize, anchor=west, inner sep=0pt, align=left] at ([xshift=\lx+2pt]axis cs:-0.28,#1)
  {\textbf{Our model: #2}\\#3};}
\def\refm#1#2{\node[font=\footnotesize, anchor=west, inner sep=0pt] at ([xshift=\lx+2pt]axis cs:-0.28,#1)
  {\hspace{0.9em}\textcolor{mute}{vs.}\hspace{0.4em}#2};}
\begin{axis}[
  scale only axis, width=5.8cm, height=7.5cm, clip=false,
  xmin=-0.28, xmax=0.045, ymin=0.3, ymax=13.1,
  axis x line*=bottom, axis y line*=left, axis line style={mute}, tick style={mute},
  y axis line style={draw=none}, ytick style={draw=none},
  xmajorgrids, grid style={black!8},
  xtick={-0.2,-0.1,0}, xticklabels={$-0.2$,$-0.1$,$0$}, minor xtick={-0.25,-0.15,-0.05}, xminorgrids, minor grid style={black!8},
  tick label style={font=\footnotesize}, label style={font=\footnotesize},
  xlabel={Loss of our model minus loss of the reference model (nats per token)},
  ytick=\empty, set layers=standard,
  legend style={font=\footnotesize, draw=none, fill=white, fill opacity=0.85, text opacity=1,
                at={(0.5,1.01)}, anchor=south, legend columns=2, /tikz/every even column/.append style={column sep=8pt}},
  legend cell align=left,
]
\pgfplotsinvokeforeach{11,9,7,5,1}{
  \fill[/pgfplots/on layer=axis background, black!5] ([xshift=\lx]axis cs:-0.28,#1-0.5) rectangle (axis cs:0.2112,#1+0.5);
}
\ours{12.3}{GDN + \ingram{}, no softmax attention}{$8\times256$, 56.0M parameters, 164M training tokens}
\refm{11}{3:1 hybrid (30.4M), same blocks and width}
\refm{10}{3:1 hybrid (54.2M), $12\times384$, parameter-matched}
\refm{9}{3:1 hybrid (55.8M) with a static $N$-gram table}
\refm{8}{3:1 hybrid (54.2M), both models $12\times384$ (ours 80.3M)}
\refm{7}{3:1 hybrid (88.4M), both $16\times512$ (ours 115.3M), 328M tokens}
\refm{6}{3:1 hybrid (30.4M), both trained for 328M tokens}
\refm{5}{3:1 hybrid (30.4M) trained twice as long, 328M tokens}
\ours{3.3}{7:1 hybrid + \ingram{}}{$8\times256$, 55.9M parameters, 1.31B training tokens}
\refm{2}{3:1 hybrid (30.4M)}
\refm{1}{7:1 hybrid (30.4M)}
\node[font=\footnotesize, text=mute, anchor=east, inner sep=0pt] at (axis cs:-0.006,12.3) {$\leftarrow$ lower loss for our model};
\draw[ink!60, line width=0.5pt] (axis cs:0,0.3) -- (axis cs:0,11.6);
\addplot[only marks, mark=*, mark size=0.85pt, cat1, opacity=0.4, forget plot]
  table[x=d, y expr=\thisrow{y}+0.17]{generated/forest_seeds.dat};
\addplot[only marks, mark=*, mark size=2.3pt, cat1,
  error bars/.cd, x dir=both, x explicit, error bar style={line width=1.1pt, cat1}, error mark options={draw=none}]
  table[x=d, y expr=\thisrow{y}+0.17, x error minus expr=\thisrow{d}-\thisrow{lo},
        x error plus expr=\thisrow{hi}-\thisrow{d}]{generated/forest.dat};
\addlegendentry{models alone (small dots: seeds)}
\addplot[only marks, mark=square*, mark size=1.9pt, cat2, mark options={fill=white, line width=0.8pt},
  error bars/.cd, x dir=both, x explicit, error bar style={line width=0.9pt, cat2}, error mark options={draw=none}]
  table[x=c, y expr=\thisrow{y}-0.2, x error minus expr=\thisrow{c}-\thisrow{clo},
        x error plus expr=\thisrow{chi}-\thisrow{c}]{generated/forest.dat};
\addlegendentry{n-gram cache in both models}
\addplot[only marks, mark=none, forget plot, point meta=explicit,
  nodes near coords={\pgfmathprintnumber[fixed, fixed zerofill, precision=3, print sign]{\pgfplotspointmeta}},
  every node near coord/.style={font=\footnotesize, text=cat1!80!black, anchor=east, inner sep=0pt}]
  table[x expr=0.1232, y expr=\thisrow{y}-0.02, meta=d]{generated/forest.dat};
\addplot[only marks, mark=none, forget plot, point meta=explicit,
  nodes near coords={\pgfmathprintnumber[fixed, fixed zerofill, precision=3, print sign]{\pgfplotspointmeta}},
  every node near coord/.style={font=\footnotesize, text=cat2!80!black, anchor=east, inner sep=0pt}]
  table[x expr=0.2025, y expr=\thisrow{y}-0.02, meta=c]{generated/forest.dat};
\node[font=\footnotesize, text=cat1!80!black, anchor=east, inner sep=0pt] at (axis cs:0.1232,12.3) {alone};
\node[font=\footnotesize, text=cat2!80!black, anchor=east, inner sep=0pt] at (axis cs:0.2025,12.3) {cache};
\end{axis}
\end{tikzpicture}
}
\caption{The pre-registered comparisons. Each bold heading names our model, and each row below it
compares that model with one reference model that has no \ingram{}. The horizontal axis is the
loss of our model minus the loss of the reference; left of zero, our model has the lower loss.
Large markers are means over
seeds 20--27 ($n = 8$; $n = 4$ in the fifth row) with 95\% bootstrap intervals, small dots are
single seeds, and open squares repeat the comparison with the n-gram cache interpolated into both
models. The columns at the right give the two means. The line under each heading gives the
blocks $\times$ width, parameters and training tokens of our model, and a row states where either
model differs; parentheses after a reference give its parameters. Both models of a row train for
the same number of steps, except in the seventh row, where the hybrid trains twice as long. Losses
and tests are in Table~\ref{tab:confirm}.}
\label{fig:confirm}
\end{figure}

Figure~\ref{fig:confirm} shows every pre-registered comparison, and Table~\ref{tab:confirm} in the
appendix gives the losses and tests in the same order. At equal training steps, a GDN with
\ingram{} and no attention reaches a lower loss than each 3:1 hybrid we trained. Against the 3:1
hybrid of the same depth and width (30.4M parameters), the loss is 0.607 against 0.749, and
$\Delta = -0.142$ nats (95\% CI $[-0.166, -0.127]$; 8 of 8 seeds; Holm $p = 0.012$). One hybrid
run is an outlier (seed 20, loss 0.835); the median difference is $-0.134$.

The model with \ingram{} has 56.0M parameters, of which 25.2M are rows of the hashed table; the
30.4M hybrid is therefore the smaller model. The second and third rows of
Figure~\ref{fig:confirm} are controls for that difference. A 3:1 hybrid with 12 blocks of width
384 has 54.2M parameters and reaches 0.709, and $\Delta = -0.102$ ($[-0.113, -0.092]$; 8 of 8;
Holm $p = 0.008$). A 3:1 hybrid that includes the same hashed table with its writes disabled, a
static $N$-gram table, has 55.8M parameters (0.1M of them are write projections that it does not
use) and reaches 0.745, against 0.749 for the hybrid
without the table, and $\Delta = -0.137$ ($[-0.179, -0.104]$; 8 of 8; Holm $p = 0.016$). Table
parameters without the writes therefore do not account for the difference. The difference also
holds at two larger sizes (fourth and fifth rows): $-0.127$ with 12 blocks of width 384 and
$-0.058$ with 16 blocks of width 512, the latter on four seeds.

The open squares of Figure~\ref{fig:confirm} interpolate the n-gram cache into both models. In
these five rows, each $\Delta$ retains its sign and decreases in magnitude to between $-0.046$ and
$-0.021$. The cache recovers 0.123 nats for the 30.4M hybrid (0.749 to 0.627) and 0.027 nats for
the model with \ingram{} (0.607 to 0.580). Against the parameter-matched hybrid the remaining
difference is $-0.023$ ($[-0.028, -0.018]$; 8 of 8; Holm $p = 0.008$).

\subsection{Longer training, and a 7:1 hybrid with \ingram{}}
\label{sec:length}

\begin{figure}[tb]
\centering
\begin{tikzpicture}
\begin{groupplot}[
  group style={group size=2 by 1, horizontal sep=1.9cm},
  width=0.46\textwidth, height=5.6cm,
  axis lines=left, axis line style={mute}, tick style={mute},
  ymajorgrids, grid style={black!8},
  tick label style={font=\footnotesize}, label style={font=\footnotesize},
  title style={font=\footnotesize},
  legend style={font=\footnotesize, draw=none, fill=white, fill opacity=0.8, text opacity=1},
  legend cell align=left,
]
\nextgroupplot[title={(a) Loss during the 1.31B-token runs},
  xlabel={Training tokens (billions)}, ylabel={Loss (nats per token)},
  xmin=0, xmax=1.4, ymin=0.4, ymax=1.3, legend pos=north east,
  yticklabel style={/pgf/number format/fixed, /pgf/number format/precision=1, /pgf/number format/zerofill}]
\addplot[name path=h7lo, draw=none, forget plot] table[x expr=\thisrow{mtok}/1000, y=h7lo]{generated/curves_p32.dat};
\addplot[name path=h7hi, draw=none, forget plot] table[x expr=\thisrow{mtok}/1000, y=h7hi]{generated/curves_p32.dat};
\addplot[cat3, opacity=0.15, forget plot] fill between[of=h7lo and h7hi];
\addplot[name path=h3lo, draw=none, forget plot] table[x expr=\thisrow{mtok}/1000, y=h3lo]{generated/curves_p32.dat};
\addplot[name path=h3hi, draw=none, forget plot] table[x expr=\thisrow{mtok}/1000, y=h3hi]{generated/curves_p32.dat};
\addplot[cat2, opacity=0.15, forget plot] fill between[of=h3lo and h3hi];
\addplot[name path=l7lo, draw=none, forget plot] table[x expr=\thisrow{mtok}/1000, y=l7lo]{generated/curves_p32.dat};
\addplot[name path=l7hi, draw=none, forget plot] table[x expr=\thisrow{mtok}/1000, y=l7hi]{generated/curves_p32.dat};
\addplot[cat1, opacity=0.15, forget plot] fill between[of=l7lo and l7hi];
\addplot[cat3, thick, mark=triangle*, mark size=2.2pt] table[x expr=\thisrow{mtok}/1000, y=h7]{generated/curves_p32.dat};
\addlegendentry{7:1 hybrid}
\addplot[cat2, thick, mark=square*, mark size=1.8pt] table[x expr=\thisrow{mtok}/1000, y=h3]{generated/curves_p32.dat};
\addlegendentry{3:1 hybrid}
\addplot[cat1, thick, mark=*, mark size=1.8pt] table[x expr=\thisrow{mtok}/1000, y=l7]{generated/curves_p32.dat};
\addlegendentry{7:1 hybrid + \ingram{}}

\nextgroupplot[title={(b) Advantage over the 3:1 hybrid within each run},
  xmode=log, ymode=log,
  xlabel={Training tokens (millions)}, ylabel={3:1 hybrid loss minus ours (nats)},
  xmin=25, xmax=1700, ymin=0.04, ymax=6,
  xtick={33,100,328,1000}, xticklabels={33,100,328,1000},
  ytick={0.05,0.1,0.2,0.5,1}, yticklabels={0.05,0.1,0.2,0.5,1},
  legend pos=north east, legend style={font=\scriptsize}]
\addplot[mute, thick, dashed, mark=o, mark size=1.8pt] table[x=mtok, y expr=-\thisrow{g}]{generated/gap_10k.dat};
\addlegendentry{GDN + \ingram{}, 164M tokens}
\addplot[mute, thick, mark=square*, mark size=1.8pt] table[x=mtok, y expr=-\thisrow{g}]{generated/gap_20k.dat};
\addlegendentry{GDN + \ingram{}, 328M tokens}
\addplot[cat1, thick, mark=*, mark size=1.8pt] table[x=mtok, y expr=-\thisrow{g13}]{generated/curves_p32.dat};
\addlegendentry{7:1 + \ingram{}, 1.31B tokens}
\end{groupplot}
\end{tikzpicture}

\caption{(a) Held-out loss during the 1.31B-token runs of the 7:1 hybrid with \ingram{}, the 3:1
hybrid and the 7:1 hybrid; lines are means over seeds 20--27 and bands span the lowest and highest
seed. The final losses are 0.448, 0.510 and 0.508 (last two rows of Figure~\ref{fig:confirm}).
(b) Loss of the 3:1 hybrid minus loss of the model with \ingram{}, at the evaluations within each
group of runs (the legend gives our model and the training tokens of the group), on logarithmic
axes. Every run follows a cosine schedule over its own length; only the last point of each curve
is a finished run.}
\label{fig:training}
\end{figure}

The difference from the 3:1 hybrid decreases when both models are trained for longer (sixth row of
Figure~\ref{fig:confirm}). With both at 20,000 steps (328M tokens), $\Delta = -0.072$ ($[-0.085,
-0.061]$; 8 of 8; Holm $p = 0.016$), half of the $-0.142$ at 164M tokens. A further pre-registered
test allows the hybrid twice the steps of our model (seventh row). It bounds a comparison at equal
wall-clock, because our training steps with \ingram{} were 1.07 to 1.19 times slower
(Appendix~\ref{sec:cost}). The model with \ingram{} retains the lower loss ($\Delta = -0.015$,
$[-0.028, -0.003]$; 6 of 8 seeds; Holm $p = 0.039$). With the n-gram cache on both sides the
hybrid has the lower loss ($+0.026$, $[+0.020, +0.032]$; 0 of 8 seeds).

The longest runs test a hybrid that retains one of the two attention layers
(Figure~\ref{fig:training}a). We trained three models for 80,000 steps (1.31B tokens, less than
one pass over the 1.65B tokens of the training set): a 7:1 hybrid with \ingram{}, the 3:1 hybrid
and the 7:1 hybrid, all with the learning rates of the 164M-token runs
(Section~\ref{sec:setup}). The 7:1 hybrid with
\ingram{} reaches 0.448, which is below the 3:1 hybrid at 0.510, the primary test ($\Delta =
-0.063$, $[-0.071, -0.054]$; 8 of 8; Holm $p = 0.008$), and below the 7:1 hybrid at 0.508
($\Delta = -0.060$, $[-0.067, -0.055]$; 8 of 8; Holm $p = 0.008$). With the n-gram cache on both
sides the differences are $-0.036$ and $-0.034$, each on 8 of 8 seeds.

Without \ingram{}, the two hybrids end at the same loss: the 3:1 hybrid minus the 7:1 hybrid is
$+0.002$ (95\% CI $[-0.011, +0.014]$; $n = 8$; post hoc). At 1.31B tokens the 7:1 hybrid with
\ingram{} is therefore 0.06 nats below a hybrid without it at either attention count. Within
these runs, its advantage over the 3:1 hybrid decreases from 0.280 nats at 8,000 steps to 0.063 at
80,000, and by 0.006 over the last 8,000 steps, when the cosine schedule is near its floor
(Figure~\ref{fig:training}b).

\subsection{Contribution of attention when \ingram{} is present}
\label{sec:types}

\begin{figure}[tb]
\centering
\begin{tikzpicture}
\begin{groupplot}[
  group style={group size=2 by 1, horizontal sep=1.9cm},
  width=0.46\textwidth, height=5.4cm,
  xlabel={Attention layers $A$}, xtick={0,1,2,4}, xmin=-0.3, xmax=4.3,
  axis lines=left, axis line style={mute}, tick style={mute},
  ymajorgrids, grid style={black!8},
  tick label style={font=\footnotesize}, label style={font=\footnotesize},
  title style={font=\footnotesize},
  yticklabel style={/pgf/number format/fixed, /pgf/number format/precision=2, /pgf/number format/zerofill},
  legend style={font=\footnotesize, draw=none, fill=white, fill opacity=0.8, text opacity=1},
  legend cell align=left,
]
\nextgroupplot[title={(a) All cells}, ylabel={Loss (nats per token)}, ymin=0.5, ymax=1.08,
  legend pos=north east]
\addplot[mute, thick, mark=o, mark size=2pt] table[x=A, y=P0]{generated/grid.dat};
\addlegendentry{$P=0$}
\addplot[ramp1, thick, mark=square*, mark size=1.8pt] table[x=A, y=P1]{generated/grid.dat};
\addlegendentry{$P=1$}
\addplot[ramp2, thick, mark=triangle*, mark size=2.2pt] table[x=A, y=P2]{generated/grid.dat};
\addlegendentry{$P=2$}
\addplot[ramp3, thick, mark=diamond*, mark size=2.2pt] table[x=A, y=P4]{generated/grid.dat};
\addlegendentry{$P=4$}
\addplot[mute, only marks, mark=o, mark size=1pt, forget plot] table[x=A, y=P0s0]{generated/grid.dat};
\addplot[mute, only marks, mark=o, mark size=1pt, forget plot] table[x=A, y=P0s1]{generated/grid.dat};

\nextgroupplot[title={(b) Cells with \ingram{}, enlarged}, ymin=0.570, ymax=0.615,
  ytick={0.57,0.58,0.59,0.60,0.61}, xmax=5.1, clip=false]
\addplot[ramp1, thick, mark=square*, mark size=1.8pt] table[x=A, y=P1]{generated/grid.dat}
  node[pos=1, right=3pt, font=\footnotesize, text=ink] {$P=1$};
\addplot[ramp2, thick, mark=triangle*, mark size=2.2pt] table[x=A, y=P2]{generated/grid.dat}
  node[pos=1, right=3pt, font=\footnotesize, text=ink] {$P=2$};
\addplot[ramp3, thick, mark=diamond*, mark size=2.2pt] table[x=A, y=P4]{generated/grid.dat}
  node[pos=1, right=3pt, font=\footnotesize, text=ink] {$P=4$};
\foreach \c/\k in {ramp1/P1, ramp2/P2, ramp3/P4} {
  \edef\tmp{\noexpand\addplot[\c, only marks, mark=o, mark size=1pt, forget plot] table[x=A, y=\k s0]{generated/grid.dat};
            \noexpand\addplot[\c, only marks, mark=o, mark size=1pt, forget plot] table[x=A, y=\k s1]{generated/grid.dat};}
  \tmp
}
\end{groupplot}
\end{tikzpicture}

\caption{Loss against the number of attention layers $A$, for $P \in \{0, 1, 2, 4\}$ placements of
\ingram{} (8 blocks of width 256, 164M training tokens). Without \ingram{} ($P = 0$) the loss
falls from 1.002 to 0.722 between no attention and two layers; with one placement all four models
lie between 0.594 and 0.609. Lines join seed means and small open circles are the two seeds. Panel
(b) repeats the models with $P \ge 1$ on a finer scale. Seeds 0 and 1 only: the grid is
exploratory. Seed means and per-seed values are in Tables~\ref{tab:grid}
and~\ref{tab:gridseeds}.}
\label{fig:grid}
\end{figure}

We trained every combination of $A \in \{0, 1, 2, 4\}$ attention layers and $P \in \{0, 1, 2, 4\}$
placements (Figure~\ref{fig:grid}). Each cell has two seeds and trains for 164M tokens; the grid
is exploratory, with its analyses fixed before it ran. Without \ingram{}, the loss depends on the
attention layers: 1.002 with none, 0.995 with one, 0.722 with two and 0.742 with four. With one
placement, all four cells lie between 0.594 and 0.609, and the largest difference between two
attention counts is 0.015 nats. Further placements lower the loss at every $A$, to between 0.575
and 0.584 at $P = 4$. One placement without any attention layer is already 0.113 nats below the
3:1 hybrid without \ingram{}.

\begin{figure}[tb]
\centering
% Generated by scripts/paper_results.py from runs/. Do not edit.
\def\atypelabels{{All tokens},{Exact repeat, far (16.4\%)},{Exact repeat, near (45.2\%)},{Consistent cue (12.1\%)},{Latest value (2.3\%)},{Older value (2.1\%)},{Novel token (5.4\%)},{No match (5.4\%)},{Near-duplicate line, far (3.7\%)},{Next integer (1.0\%)},{Format (4.7\%)}}

\begin{tikzpicture}
\begin{axis}[
  scale only axis, width=8.6cm, height=6.2cm,
  xmin=-0.04, xmax=0.82, ymin=0.4, ymax=11.6,
  axis x line*=bottom, axis y line*=left, axis line style={mute}, tick style={mute},
  y axis line style={draw=none}, ytick style={draw=none},
  xmajorgrids, grid style={black!8},
  xtick={0,0.2,0.4,0.6,0.8},
  xticklabel style={/pgf/number format/fixed, /pgf/number format/precision=1, /pgf/number format/zerofill},
  tick label style={font=\footnotesize}, label style={font=\footnotesize},
  xlabel={Loss without attention minus loss with two attention layers (nats per token)},
  ytick={11,10,9,8,7,6,5,4,3,2,1},
  yticklabels/.expand once=\atypelabels,
  yticklabel style={font=\footnotesize, align=right},
  legend style={font=\footnotesize, draw=none, fill=white, fill opacity=0.85, text opacity=1,
                at={(0.98,0.03)}, anchor=south east},
  legend cell align=left,
]
\draw[ink!60, line width=0.5pt] (axis cs:0,0.4) -- (axis cs:0,11.6);
\draw[black!15, line width=0.4pt] (axis cs:-0.04,10.5) -- (axis cs:0.82,10.5);
\addplot[only marks, mark=o, mark size=2.3pt, ink!70, mark options={line width=0.9pt},
  error bars/.cd, x dir=both, x explicit, error bar style={line width=0.9pt, ink!55}, error mark options={draw=none}]
  table[x=c0, y expr=\thisrow{y}+0.15, x error minus expr=\thisrow{c0}-\thisrow{lo0},
        x error plus expr=\thisrow{hi0}-\thisrow{c0}]{generated/a3.dat};
\addlegendentry{without \ingram{}}
\addplot[only marks, mark=*, mark size=2.3pt, cat1,
  error bars/.cd, x dir=both, x explicit, error bar style={line width=1.1pt, cat1}, error mark options={draw=none}]
  table[x=c1, y expr=\thisrow{y}-0.15, x error minus expr=\thisrow{c1}-\thisrow{lo1},
        x error plus expr=\thisrow{hi1}-\thisrow{c1}]{generated/a3.dat};
\addlegendentry{with \ingram{}, one placement}
\end{axis}
\end{tikzpicture}

\caption{Cost of removing both attention layers of a 3:1 hybrid, by token type: the loss with no
attention layer minus the loss with two, without \ingram{} (open circles) and with one placement
(filled). Without \ingram{} the cost is largest on far exact repeats (0.646 nats); with one
placement it is 0.000 there and 0.015 over all tokens. Bars are 95\% bootstrap intervals over the
32 evaluation windows, with seeds 0 and 1 averaged (exploratory). Percentages are shares of
evaluation tokens; types overlap and are defined in Table~\ref{tab:tags}. Table~\ref{tab:a3} gives
all placements and the one type with under 0.5\% of tokens.}
\label{fig:types}
\end{figure}

On the same grid, we decompose the cost of removing both attention layers, $C(P) = L(A{=}0, P) -
L(A{=}2, P)$, by token type (Figure~\ref{fig:types}). Without \ingram{} it is
0.280 nats per token ($[0.215, 0.344]$). Among the types with at least 1\% of tokens, it is
largest on far exact repeats, the tokens that continue a match of at least 8 tokens whose previous
occurrence lies more than 512 tokens back ($+0.646$), and on cues whose earlier continuations all
agree ($+0.465$). With one placement the cost over all tokens is 0.015 ($[0.010, 0.020]$) and the
cost on far exact repeats is 0.000 ($[-0.003, 0.003]$).

The remaining cost at $P = 1$ is concentrated on three types, with 2.1\%, 5.4\% and 4.7\% of
tokens. These are older-value tokens, where the matched cue had several earlier continuations and
the most recent one is wrong ($+0.053$, $[0.036, 0.073]$); tokens with no earlier match ($+0.053$,
$[0.026, 0.099]$); and format tokens ($+0.041$, $[0.013, 0.102]$). In the confirmatory runs
(Table~\ref{tab:types}; post hoc), far exact repeats are 16.4\% of tokens and likewise account for 60\%,
76\% and 57\% of the difference from the 3:1 hybrid at 164M, 328M and 1.31B tokens. Without
attention, the model with \ingram{} is worse than the 3:1 hybrid on older-value tokens ($+0.089$,
$[0.042, 0.138]$ at 328M tokens). With one attention layer, at 1.31B tokens, this type no longer
favours the 3:1 hybrid ($-0.013$, $[-0.027, 0.000]$).
